% related_work.tex (chundawat-h)

\section{Related work}
\label{sec:related}

\paragraph{Receipt and form understanding.} The ICDAR 2019 Robust
Reading Challenge on Scanned Receipts OCR and Information Extraction
(SROIE)~\cite{huang2019sroie} introduced a public benchmark of about
1{,}000 scanned receipts with three tasks: text localisation, OCR, and
key information extraction of company, date, address and total.
CORD~\cite{park2019cord} provides a larger receipt dataset with more
fine-grained labels. Before Transformers, such fields were usually
extracted with hand-crafted rules and templates, which work well for
fixed layouts but break on unseen ones.

\paragraph{Sequence labelling.} Key information extraction is commonly
framed as named-entity-style token classification with BIO
tags~\cite{ramshaw1995text}: each word is labelled as the Beginning or
Inside of a field, or Outside any field. Neural taggers such as the
BiLSTM-CRF~\cite{lample2016ner} became the standard approach and were
later replaced by fine-tuned Transformers~\cite{vaswani2017attention}.

\paragraph{Pre-trained language models.} BERT~\cite{devlin2019bert} is a
Transformer encoder pre-trained on large text corpora with masked
language modelling. Adding a classification layer and fine-tuning on
labelled tokens gives strong results on entity recognition. However,
BERT reads a document as a 1D sequence and cannot see where a word is on
the page.

\paragraph{Layout-aware models.} Chargrid~\cite{katti2018chargrid}
encodes a document as a 2D grid of characters for a convolutional
network. LayoutLM~\cite{xu2020layoutlm} keeps the BERT architecture and
adds 2D position embeddings of each word's bounding box. It was
pre-trained on 11 million scanned document pages and reported large
gains over BERT on form understanding and on SROIE.
LayoutLMv2~\cite{xu2021layoutlmv2} further adds visual features from the
page image. We use LayoutLM (v1) because it differs from BERT
\emph{only} in the layout input, which makes the comparison clean.
